%
\documentclass[12pt]{article}

\usepackage[usenames]{color}
\usepackage{verbatim}
\usepackage{fancyhdr}
\usepackage{graphicx}
\usepackage{amssymb,amsthm}
%\usepackage[fleqn]{amsmath}
\usepackage{amsmath}
% \usepackage[final,notref,notcite,color]{showkeys}
\usepackage[left=1in,top=1in,right=1in,bottom=1in,letterpaper]{geometry}
\renewcommand{\baselinestretch}{1.2}

\pagestyle{fancy}
\lhead{}
\chead{}
\rhead{}
\lfoot{}
\cfoot{\thepage}
\rfoot{}



%% editing
\usepackage[shortlabels]{enumitem}
\usepackage[normalem]{ulem}

%% macros for commenting
\usepackage{mathtools}
\usepackage[normalem]{ulem} % to use \sout

%% math packages
% \usepackage{amsmath} % not needed for Springer
\usepackage{amssymb}
\usepackage{mathrsfs}

%\usepackage{subfigure}
%\usepackage{url}

%% a light-weight algorithm environment
\newtheorem{algo}{Algorithm}

%% highlighting and commenting
\newcommand{\outline}[1]{{\color{brown}#1}}
\newcommand{\rev}[1]{{\color{blue}#1}}
\newcommand{\remove}[1]{{\sout{#1}}}
\newcommand{\cut}[1]{{}}

%% macros for letters

\newcommand{\va}{{\mathbf{a}}}
\newcommand{\vb}{{\mathbf{b}}}
\newcommand{\vc}{{\mathbf{c}}}
\newcommand{\vd}{{\mathbf{d}}}
\newcommand{\ve}{{\mathbf{e}}}
\newcommand{\vf}{{\mathbf{f}}}
\newcommand{\vg}{{\mathbf{g}}}
\newcommand{\vh}{{\mathbf{h}}}
\newcommand{\vi}{{\mathbf{i}}}
\newcommand{\vj}{{\mathbf{j}}}
\newcommand{\vk}{{\mathbf{k}}}
\newcommand{\vl}{{\mathbf{l}}}
\newcommand{\vm}{{\mathbf{m}}}
\newcommand{\vn}{{\mathbf{n}}}
\newcommand{\vo}{{\mathbf{o}}}
\newcommand{\vp}{{\mathbf{p}}}
\newcommand{\vq}{{\mathbf{q}}}
\newcommand{\vr}{{\mathbf{r}}}
\newcommand{\vs}{{\mathbf{s}}}
\newcommand{\vt}{{\mathbf{t}}}
\newcommand{\vu}{{\mathbf{u}}}
\newcommand{\vv}{{\mathbf{v}}}
\newcommand{\vw}{{\mathbf{w}}}
\newcommand{\vx}{{\mathbf{x}}}
\newcommand{\vy}{{\mathbf{y}}}
\newcommand{\vz}{{\mathbf{z}}}

\newcommand{\ta}{{\tilde{a}}}
\newcommand{\tb}{{\tilde{b}}}
\newcommand{\tc}{{\tilde{c}}}
\newcommand{\td}{{\tilde{d}}}
\newcommand{\te}{{\tilde{e}}}
\newcommand{\tf}{{\tilde{f}}}
\newcommand{\tg}{{\tilde{g}}}

\newcommand{\ti}{{\tilde{i}}}
\newcommand{\tj}{{\tilde{j}}}
\newcommand{\tk}{{\tilde{k}}}
\newcommand{\tl}{{\tilde{l}}}
\newcommand{\tm}{{\tilde{m}}}
\newcommand{\tn}{{\tilde{n}}}

\newcommand{\tp}{{\tilde{p}}}
\newcommand{\tq}{{\tilde{q}}}

\newcommand{\ts}{{\tilde{s}}}

\newcommand{\tu}{{\tilde{u}}}
\newcommand{\tv}{{\tilde{v}}}
\newcommand{\tw}{{\tilde{w}}}
\newcommand{\tx}{{\tilde{x}}}
\newcommand{\ty}{{\tilde{y}}}
\newcommand{\tz}{{\tilde{z}}}

\newcommand{\vA}{{\mathbf{A}}}
\newcommand{\vB}{{\mathbf{B}}}
\newcommand{\vC}{{\mathbf{C}}}
\newcommand{\vD}{{\mathbf{D}}}
\newcommand{\vE}{{\mathbf{E}}}
\newcommand{\vF}{{\mathbf{F}}}
\newcommand{\vG}{{\mathbf{G}}}
\newcommand{\vH}{{\mathbf{H}}}
\newcommand{\vI}{{\mathbf{I}}}
\newcommand{\vJ}{{\mathbf{J}}}
\newcommand{\vK}{{\mathbf{K}}}
\newcommand{\vL}{{\mathbf{L}}}
\newcommand{\vM}{{\mathbf{M}}}
\newcommand{\vN}{{\mathbf{N}}}
\newcommand{\vO}{{\mathbf{O}}}
\newcommand{\vP}{{\mathbf{P}}}
\newcommand{\vQ}{{\mathbf{Q}}}
\newcommand{\vR}{{\mathbf{R}}}
\newcommand{\vS}{{\mathbf{S}}}
\newcommand{\vT}{{\mathbf{T}}}
\newcommand{\vU}{{\mathbf{U}}}
\newcommand{\vV}{{\mathbf{V}}}
\newcommand{\vW}{{\mathbf{W}}}
\newcommand{\vX}{{\mathbf{X}}}
\newcommand{\vY}{{\mathbf{Y}}}
\newcommand{\vZ}{{\mathbf{Z}}}

\newcommand{\cA}{{\mathcal{A}}}
\newcommand{\cB}{{\mathcal{B}}}
\newcommand{\cC}{{\mathcal{C}}}
\newcommand{\cD}{{\mathcal{D}}}
\newcommand{\cE}{{\mathcal{E}}}
\newcommand{\cF}{{\mathcal{F}}}
\newcommand{\cG}{{\mathcal{G}}}
\newcommand{\cH}{{\mathcal{H}}}
\newcommand{\cI}{{\mathcal{I}}}
\newcommand{\cJ}{{\mathcal{J}}}
\newcommand{\cK}{{\mathcal{K}}}
\newcommand{\cL}{{\mathcal{L}}}
\newcommand{\cM}{{\mathcal{M}}}
\newcommand{\cN}{{\mathcal{N}}}
\newcommand{\cO}{{\mathcal{O}}}
\newcommand{\cP}{{\mathcal{P}}}
\newcommand{\cQ}{{\mathcal{Q}}}
\newcommand{\cR}{{\mathcal{R}}}
\newcommand{\cS}{{\mathcal{S}}}
\newcommand{\cT}{{\mathcal{T}}}
\newcommand{\cU}{{\mathcal{U}}}
\newcommand{\cV}{{\mathcal{V}}}
\newcommand{\cW}{{\mathcal{W}}}
\newcommand{\cX}{{\mathcal{X}}}
\newcommand{\cY}{{\mathcal{Y}}}
\newcommand{\cZ}{{\mathcal{Z}}}

\newcommand{\ri}{{\mathrm{i}}}
\newcommand{\rr}{{\mathrm{r}}}

\newcommand{\EE}{{\mathbb{E}}}


%% macros for math notions and operators

\newcommand{\RR}{\mathbb{R}}
\newcommand{\CC}{\mathbb{C}}
\newcommand{\ZZ}{\mathbb{Z}}
\renewcommand{\SS}{{\mathbb{S}}}
\newcommand{\SSp}{\mathbb{S}_{+}}
\newcommand{\SSpp}{\mathbb{S}_{++}}
\newcommand{\sign}{\mathrm{sign}}
\newcommand{\vzero}{\mathbf{0}}
\newcommand{\vone}{{\mathbf{1}}}

\newcommand{\st}{{\text{s.t.}}} % subject to
\newcommand{\St}{{\mathrm{subject~to}}} % subject to
\newcommand{\op}{{\mathrm{op}}} % subscript for operator norm
\newcommand{\opt}{{\mathrm{opt}}} % subscript for optimal solution
%\newcommand{\supp}{{\mathrm{supp}}} % support
\newcommand{\Prob}{{\mathrm{Prob}}} % probability
\newcommand{\Diag}{{\mathrm{Diag}}} % vector -> diagonal matrix
%\newcommand{\diag}{{\mathrm{diag}}} % matrix diagonal -> vector
\newcommand{\dom}{{\mathrm{dom}}} % domain
\newcommand{\range}{{\mathrm{range}}} % domain
%\newcommand{\grad}{{\nabla}}    % gradient
\newcommand{\tr}{{\mathrm{tr}}} % trace
\newcommand{\TV}{{\mathrm{TV}}} % total variation
\newcommand{\Proj}{{\mathrm{Proj}}}
\newcommand{\prj}{{\mathrm{prj}}}
\newcommand{\prox}{\mathbf{prox}}
\newcommand{\refl}{\mathbf{refl}}
\newcommand{\reflh}{\refl^{\bH}}
\newcommand{\proxh}{\prox^{\bH}}
\newcommand{\minimize}{\text{minimize}}
\newcommand{\bgamma}{\boldsymbol{\gamma}}
\newcommand{\bsigma}{\boldsymbol{\sigma}}
\newcommand{\bomega}{\boldsymbol{\omega}}
\newcommand{\blambda}{\boldsymbol{\lambda}}
\newcommand{\bH}{\vH}
\newcommand{\bbH}{\mathbb{H}}
\newcommand{\bB}{\boldsymbol{\cB}}
\newcommand{\Tau}{\mathrm{T}}
\newcommand{\tnabla}{\widetilde{\nabla}}
\newcommand{\TDRS}{T_{\mathrm{DRS}}}
\newcommand{\TPRS}{T_{\mathrm{PRS}}}
\newcommand{\TFBS}{T_{\mathrm{FBS}}}
\newcommand{\best}{\mathrm{best}}
\newcommand{\kbest}{k_{\best}}
\newcommand{\diff}{\mathrm{diff}}
\newcommand{\barx}{\bar{x}}
\newcommand{\xgbar}{\bar{x}_g}
\newcommand{\xfbar}{\bar{x}_f}
\newcommand{\hatxi}{\hat{\xi}}
\newcommand{\xg}{x_g}
\newcommand{\xf}{x_f}
\newcommand{\du}{\mathrm{d}u}
\newcommand{\dy}{\mathrm{d}y}
\newcommand{\kconvergence}{\stackrel{k \rightarrow \infty}{\rightarrow }}
\DeclareMathOperator{\shrink}{shrink} % shrinkage
\DeclareMathOperator*{\argmin}{arg\,min}
\DeclareMathOperator*{\argmax}{arg\,max}
\DeclareMathOperator*{\Min}{minimize}
\DeclareMathOperator*{\Max}{maximize}
\DeclareMathOperator*{\Fix}{Fix}
\DeclareMathOperator*{\zer}{zer}
\DeclareMathOperator*{\nablah}{\nabla^{\bH}}
\DeclareMathOperator*{\gra}{gra}
\DeclarePairedDelimiter{\dotpb}{\langle}{\rangle_{\bH}}
\DeclarePairedDelimiter{\dotpv}{\langle}{\rangle_{\vH}}
\DeclarePairedDelimiter{\dotp}{\langle}{\rangle}

%% macros for environments math equations

\newcommand{\MyFigure}[1]{../fig/#1}

\newcommand{\bc}{\begin{center}}
\newcommand{\ec}{\end{center}}

\newcommand{\bdm}{\begin{displaymath}}
\newcommand{\edm}{\end{displaymath}}

\newcommand{\beq}{\begin{equation}}
\newcommand{\eeq}{\end{equation}}

\newcommand{\bfl}{\begin{flushleft}}
\newcommand{\efl}{\end{flushleft}}

\newcommand{\bt}{\begin{tabbing}}
\newcommand{\et}{\end{tabbing}}

\newcommand{\beqn}{\begin{align}}
\newcommand{\eeqn}{\end{align}}

\newcommand{\beqs}{\begin{align*}} % no equation numbers
\newcommand{\eeqs}{\end{align*}}  % no equation numbers

%% macros for theorem-like environments

%\newtheorem{theorem}{Theorem}
\newtheorem{assumption}{Assumption}
%\newtheorem{condition}{Condition}
%\newtheorem{rul}{Rule}
%\newtheorem{definition}{Definition}
%\newtheorem{corollary}{Corollary}
%\newtheorem{remark}{Remark}
%\newtheorem{lemma}{Lemma}
%\newtheorem{proposition}{Proposition}
%\newtheorem{example}{Example}
%\newtheorem{proof}{Proof}

%%%%%%%%%%%%%%%%%%%%%%%%%%%%%%%%%%%%%%%%%%%%%

\title{Assignment 7 of MATP4820/6610 }
\author{(Due at 11:59pm on April-09-2026 in GradeScope)}
\date{}



\begin{document}

\maketitle

\noindent\textbf{Requirement}: you must write the solutions by yourself. As different problems may be assigned to MATP4820 and 6610 students, \textbf{you must explicitly write whether you are taking MATP4820 or 6610}. 

\vspace{0.2cm}

\noindent\textbf{Bonus}: In addition to the bonus problems (that will be assigned occasionally), you can earn up to 5\% bonus credit if you write your homework in LaTex. 

\section*{Problem 1}
Consider the problem
$$\min_{x,y\in\RR} f(x,y) = \frac{1}{2}x^2-xy+4y^2+5x-12y,\ \st\ x\le 2, y\le 4.$$

\begin{enumerate}
\item Does $(0,1)$ satisfy the optimality condition?
\item Show that $f$ is convex on $\RR^2$. Find the global minimizer.
\end{enumerate}

\newpage

\section*{Problem 2}
Consider the problem
\begin{align*}
\min_{x,y\in\RR} &~f(x,y) = -x^2+xy+\frac{1}{2}y^2-2x-y\\
\st &~ x\ge -1, y \ge -1, x+y\le 4.
\end{align*}
Find all KKT points and the global minimizer.

\newpage


\section*{Problem 3}
Consider the nonnegative quadratic program:
\begin{equation}\label{eq:nqp}
\Min_{\vx\in X} f(\vx) = \frac{1}{2}\vx^\top \vQ\vx - \vc^\top\vx, \st\ \vA\vx = \vb
\end{equation}
where $\vQ$ is a symmetric and positive semidefinite matrix, and $X=\{\vx\in\RR^n: x_i\ge0, i=1,\ldots, n\}$.

Use the instructor's provided file \verb|penalty_qp.m| to write a Matlab function \verb|penalty_qp| with input $\vQ, \vc, \vA$, $\vb$, initial vector $\vx0$, stopping tolerance \verb|tol|, initial penalty parameter $\mu_0$, and the final $\mu_1$. Also test your function by running the provided test file \verb|test_penalty_qp.m| and compare to the instructor's function. Print your code and the results you get. [\textbf{Do your best to optimize the efficiency of your code. You will lose 10\% marks if your solver takes more than double of time of the instructor's solver, and lose 20\% if your solver takes more than triple of the instructor's solver.}]

\newpage

\section*{Problem 4}
Consider the one-dimensional minimization problem
\begin{equation}\label{eq:prob1}
\min_x 2x - 1,\ \st \ \frac{1}{2}x\le 1,\, x+1 \ge 0
\end{equation}
\begin{enumerate}
\item Formulate the quadratic penalty problem with a penalty parameter $\mu > 0$ by penalizing \textbf{both} inequality constraints.
\item Derive the optimal solution of the quadratic penalty problem. Note that the solution should depend on $\mu$.
\item Prove that the solution of the quadratic penalty problem converges to the optimal solution $x^*=-1$ of Problem \eqref{eq:prob1}.
\end{enumerate}

\newpage

\section*{Problem 5}
Consider the problem
$$\min_{x,y\in\RR} f(x,y) = -\frac{1}{2}x^2-xy + \frac{1}{2}y^2-2x-2y,\ \st\ x+y=1.$$

\begin{enumerate}
\item Argue that $(1,0)$ is the unique minimizer.
\item The augmented Lagrangian function with penalty parameter $\beta>0$ is
$$\cL(x,y,v) = -\frac{1}{2}x^2-xy + \frac{1}{2}y^2-2x-2y + v(x+y-1)+ \frac{\beta}{2}(x+y-1)^2.$$
Argue that for any $\beta\in(0,1]$, the augmented Lagrangian method with initial multiplier $v^{(0)}=0$ does not work for this instance, i.e., the generated sequence $(x^{(k)},y^{(k)})$ does not converge to a KKT point.
\item For $\beta = 2$, do by hand two iterations of the augmented Lagrangian method with initial multiplier $v^{(0)}=0$.
\item (\textbf{bonus question}) For any $\beta>2$, does the generated sequence $(x^{(k)},y^{(k)})$ converge to the unique minimizer $(1,0)$, starting from any $v^{(0)}$? If yes, prove it; if not, give a value of $\beta>2$ and $v^{(0)}$ such that the sequence does not converge to $(1,0)$.
\end{enumerate}

\end{document}
